\documentclass{article}

\usepackage[preprint]{cs224r_2026}
\usepackage[utf8]{inputenc}
\usepackage[T1]{fontenc}
\usepackage[hidelinks]{hyperref}
\usepackage{url}
\usepackage{booktabs}
\usepackage{amsfonts}
\usepackage{nicefrac}
\usepackage{microtype}
\usepackage{xcolor}
\usepackage{amsmath}
\usepackage{fancyhdr}
\pagestyle{fancy}
\fancyhf{}
\renewcommand{\headrulewidth}{0pt}

\renewcommand{\maketitle}{
  \begin{center}
    {\bf \large Enhancing SINGER: Onboard Visual Drone Navigation through Iterative Imitation Learning and Reinforcement Fine-Tuning}\\
    \vspace{0.2cm}
    {\bf Authors:} Erwin Poussi, Maximilian Adang, Mac Schwager\\
    \vspace{0.1cm}
    Stanford Multi-Robot Systems Lab\\
    \vspace{0.3cm}
  \end{center}
}

\begin{document}

\maketitle

\section*{Abstract}

Autonomous drone navigation from onboard vision and natural language commands remains a challenging problem because the policy must simultaneously localize a semantically specified target, take collision-free actions, and output low-level rate commands, all from a single RGB stream without access to depth sensors or ground-truth positions. SINGER proposes a language-conditioned policy learned through behavioral cloning on expert RRT trajectories in a photorealistic 3D Gaussian Splatting environment, but yet still presents a gap to bridge since a comprehensive SINGER training achieves roughly 70\% success with frequent collisions. We propose a pilot study of iterative DAgger enhancement of SINGER that bridges the distributional shift and multi-modal averaging of the BC-learned policy, combined with an improvement of the input representation by inferring visual goal-position features from CLIPSeg semantic segmentation, namely bearing and elevation angles to the target object.

We perform our study on a reduced pilot training dataset of approximately 150,000 samples compared to the 1.5 million used in SINGER's canonical training. With our settings the SINGER baseline BC training achieves 36.7\% success rate. When adding DAgger alone success rises to 52\%, when combining BC with visually inferred goal-position features up to 80\% with the exact same 150,000 training samples, and when pairing DAgger with visually inferred goal-position features up to 90\% on average with sub-8\% collision rate. These results demonstrate both the efficacy of DAgger but also the criticality of goal-position inference. Failure modes are mainly related to goal detection problems, which paves the way for improving reliable detection in future work and for potential reinforcement learning fine-tuning to overcome noisy spurious expert-gathered data.

\end{document}
